\section{introduction}

One of the most important goals of quantum chromodynamics (QCD) is to understand the hadron structure from its fundamental degrees of freedom --- quarks and gluons.
A powerful tool of obtaining such results is lattice QCD, which has been used to study the static properties of the hadron, such as mass, charge radius etc., to good accuracy~(see e.g.~\cite{Edwards:2012fx,Aoki:2016frl}). However, in general it cannot be used to directly access quantities with real-time dependence. An example is the parton distribution functions (PDFs), which are defined as the quark and gluon matrix elements of light-cone correlations and describe the momentum distribution of quarks and gluons inside the hadron. They play a crucial role in understanding the experimental data at high energy hadron colliders such as the LHC. Lattice QCD can only be used to indirectly extract the information of the PDFs by calculating their moments~\cite{Detmold:2001jb,Dolgov:2002zm}, which is limited by technical difficulties.

In the past few years, a new approach has been proposed to study parton physics from lattice QCD, which is now known as the large-momentum effective theory (LaMET)~\cite{Ji:2013dva,Ji:2014gla}. Here one starts from a time-independent quasi-PDF, which is a matrix element of non-local quark and gluon operators that can be directly calculated on the lattice, and then uses it to extract the light-cone PDF through a perturbative factorization up to power corrections suppressed by the nucleon momentum. The factorization in Refs.~\cite{Ji:2013dva,Xiong:2013bka} was given for the bare quark quasi-PDF, where all fields and couplings entering the quasi-distribution are bare ones.
The ultraviolet (UV) physics of the quasi-PDF is all factorized perturbatively into the matching coefficients. However, lattice perturbation theory is known to converge slowly, and non-perturbative renormalization is often necessary to define a continuum limit where the matching to
parton physics is under better control. Thus to make LaMET more practical, a complete understanding of UV divergences of the relevant non-local operators
is essential.

Renormalization of non-local operators in QCD has not been well studied in the literature. A distinct feature of the LaMET non-local operators is the appearance of Wilson lines connecting the parton fields at different spacetime points to ensure gauge invariance.
The renormalization of an open Wilson line as well as a closed Wilson loop has been studied decades ago~\cite{Dotsenko:1979wb,Craigie:1980qs,Dorn:1986dt}, where it was shown that the Wilson line with a smooth simple contour can be renormalized, and in particular, the power divergence associated with the Wilson line can be absorbed into an exponential mass renormalization. In an auxiliary field formalism~\cite{Dorn:1986dt}, it was also shown that the Wilson line can be replaced by a Green's function of the auxiliary field and studied as such. On the other hand, the renormalization of non-local operators such as the quasi-PDF has extra complications due to the parton fields attached to the endpoints.
We have shown previously that the quasi-PDF is multiplicatively renormalizable in dimensional regularization up to two-loop order~\cite{Ji:2015jwa}.
The full renormalization properties of the quasi-PDF has been conjectured in recent studies~\cite{Ishikawa:2016znu,Chen:2016fxx,Constantinou:2017sej,Alexandrou:2017huk,Chen:2017mzz,Stewart:2017ss} (for an earlier work on the renormalization of transverse-momentum-dependent distributions with a similar conjecture, see~\cite{Musch:2010ka}),
but not proven.

In this paper, we perform a systematic study of the renormalization of gauge-invariant non-local operators for the quasi-PDF. We first present a general framework by introducing an auxiliary ``heavy quark'' field into the QCD Lagrangian, where the Wilson line can be replaced by a product of the auxiliary fields. In analogy to heavy-quark effective theory (HQET), we argue that this theory is renormalizable to all orders in perturbation theory. By taking the unpolarized quark quasi-PDF as an example, we then show that its renormalization reduces to that of two heavy-light quark currents, which can also be done to all orders in perturbation theory. This removes
the obstacle to define a continuous quasi-PDF through lattice simulations. The same conclusion is not limited to the quark quasi-PDF, but can also be generalized to the gluon ones, as well as
other non-local correlators, such as the Euclidean observables defined in the LaMET framework to extract the generalized parton distributions,
transverse-momentum-dependent distributions, hadron distribution
amplitudes, etc.
